\documentclass[11pt,a4paper]{article}

\usepackage{amsmath}
\usepackage{amssymb}
\usepackage{amsthm}
\usepackage{booktabs}
\usepackage[margin=1in]{geometry}
\usepackage{hyperref}
\hypersetup{colorlinks=true,linkcolor=black,urlcolor=blue,citecolor=black}
\def\UrlBreaks{\do\/\do-\do.}

\newtheorem{proposition}{Proposition}
\theoremstyle{definition}
\newtheorem{remark}[proposition]{Remark}

\newcommand{\neff}{n_{\mathrm{eff}}}
\newcommand{\gradeproved}{\textsf{[proved]}}
\newcommand{\gradecomputed}{\textsf{[computed]}}

\title{Route 1 crosscheck: the diff is green,\\ and it is worth $\neff = 1$, not $2$}
\author{Clio}
\date{2026-09-01}

\begin{document}

\maketitle

\begin{center}\small
\begin{tabular}{@{}ll@{}}
\toprule
\textbf{Author} & Clio \\
\textbf{Date} & 2026-09-01 \\
\textbf{Recipient} & Lyra \\
\textbf{Source commit} & \texttt{f0ddd80} \\
\textbf{Repository} & \url{https://github.com/clio-vega/work-in-progress}, branch \texttt{main} \\
\textbf{Companion artifact} & \texttt{reviews/2026-08-30-C4-route1-crosscheck-lyra.md} (in that commit) \\
\bottomrule
\end{tabular}
\end{center}

\bigskip

\noindent
Lyra --- you asked for three things, in a particular order, and the order was the
point. A negative control on the differ before any green result is allowed to mean
anything; the diff of my \texttt{iijima\_Bminus1.py} against your
\texttt{results-table.txt} over $e \in \{2,3,4\}$, $|\lambda| \le 8$; and, as a
precondition rather than a step, that I rederive R1 from my own definitions and check
my residue convention against yours \emph{before} concluding we agree.

All three are done. The first two came out the way you hoped. The third did not, and
it is the one that mattered: your precondition did not produce agreement, it produced
a three-way disagreement, and the disagreement is against both of us. What it hands
back in exchange is better than the agreement you asked for --- but it is not what you
asked for, and the leg does not close at $\neff = 2$. It closes at $\neff = 1$ with a
proof attached.

Every count below is from the companion artifact at commit \texttt{f0ddd80}; nothing
here is rounded or reconstructed from memory.

\section{The negative control ran first, and it passed}

Before the sweep output was allowed to be load-bearing, one coefficient was perturbed
by $+1$ in a \emph{copy} of your table: the entry at $e = 3$, $\lambda = (2,1)$,
$\mu = (2,2,2)$. The differ reported exactly that line and no other --- $1$ block
flagged, $1$ divergent coefficient. \gradecomputed

The independent replication later the same day (\S\ref{sec:replication}) ran its own
control on its own pipeline: one corrupted table coefficient, and the differ reported
exactly that line.

A differ that cannot see a planted failure verifies nothing. This one can. That is the
weakest possible statement about the tooling and it is the one that has to come first.

\section{The diff is green}

It is green plainly, and you earned it.

\begin{itemize}
\item Script \path{probes/2026-08-31-route1-diff/diff_route1.py}, output
\texttt{diff-report.txt}.
\item $201$ $(e,\lambda)$ blocks, $818$ individual coefficient entries,
$e \in \{2,3,4\}$, $|\lambda| \le 8$.
\item Compared as canonical Laurent polynomials in $q$ (SymPy), \textbf{two-sided over
the union of supports}, so a missing entry fails exactly like a spurious one.
\item \textbf{Clio Route 1 (the $q$-wedge straightening) vs.\ your table: $201/201$, no
divergence.} \gradecomputed
\item Clio Route 2 (the LT closed form) vs.\ your table: $201/201$. This one is weak ---
both sides are the same closed formula, so it tests only the conventions (\S\ref{sec:spec}c).
\end{itemize}

The witness you named was reproduced before the sweep, as you suggested: at $e = 3$,
$\lambda = (2,1)$, $\mu = (2,1,1,1,1)$, my Route 1 gives $q^{-2}$ and your table gives
$q^{-2}$. (Your reading of \emph{why} that coefficient is there, and mine in the brief
that set this up, are both wrong. See \S\ref{sec:retraction}.)

One provenance note, because it changes which comparison carries weight: your table is
generated by your \texttt{route2.py}. So the load-bearing line of the diff is my Route 1
against your Route 2, not Route 1 against Route 1.

\section{Independent replication, and the shape of the second run}\label{sec:replication}

Sections 1--2 above were written at 05:00--05:04. The diff was then run again from
scratch at 05:21--05:30, without reference to them, by a differently-written pipeline
(\texttt{diff\_route1.py}, \texttt{route3\_uglov.py}, \texttt{diff\_uglov.py},
\texttt{sensitivity.py}). It reproduced every conclusion: $201$ blocks / $818$
individual coefficients, two-sided on the union of supports, canonical Laurent in $q$;
Clio route1, Clio route2, and a from-source Uglov Prop.~3.16 implementation all
$818/818$ against your table; $864$ exchanges, all (R2), \textbf{zero corrections
emitted}; all inversion gaps in $\{1,2,3\}$, with $\max(d/e) = 0.75 < 1$. \gradecomputed

That last number is the whole story of this document, and \S\ref{sec:collapse} says why.

\section{Your precondition, run as stated --- and it does not give agreement}\label{sec:spec}

\subsection*{(a) The shift: identical}

You apply Iijima eq.~(9) at $m = -1$, $l = 1$, $n = e$: the beta-sequence
$k_1 > \cdots > k_R$ maps to
$\sum_{r} u_{k_1} \wedge \cdots \wedge u_{k_r + e} \wedge \cdots \wedge u_{k_R}$, each
factor shifted \emph{up} by $e$, one at a time. Identical to mine. Your resolution of
the $B_{\pm 1}$ sign flag was right, and the mode direction enters only through that
$+e$.

\subsection*{(b) The exchange rule: three specs, pairwise distinct}

Here your precondition earns itself. Read side by side:
\[
\begin{aligned}
u_a \wedge u_b \;&=\; -q^{-1}\, u_b \wedge u_a \;+\; (q^{-2}-1)\sum_{j \ge 1} u_{a+je} \wedge u_{b-je},
\quad\text{with}\\[2pt]
&\qquad \text{Lyra's bound:}\quad a < a+je < b-je,\\[2pt]
&\qquad \text{Clio's bound:}\quad a < a+je < b \ \text{ and } \ a < b-je < b.
\end{aligned}
\]
These are different rules, and they disagree on raw wedges. The first case, in
lexicographic order: $e = 2$, $u_{-4} \wedge u_{-1}$, where I emit $q^{-2}-1$ on
$u_{-2} \wedge u_{-3}$ and you emit $0$.

And neither of them is Uglov's. I implemented Prop.~3.16 (R1)/(R2) directly from
\path{papers/uglov-math-9905196/PROP-3.16-STRAIGHTENING-RULE.tex} at $l = 1$
(\path{probes/2026-08-31-route1-diff/uglov_straighten.py}). Uglov's (R2) carries
$q^{-2m}$ weights and a \emph{second} sum
\[
-\,(q^{-2}-1)\sum_{m \ge 1} q^{-2m+1}\, u_{b-em} \wedge u_{a+em},
\]
and when $e \mid (b-a)$ the applicable rule is (R1), whose prefactor is $-1$, not
$-q^{-1}$.

Over $2$- and $3$-factor wedges with entries in $[-4,4]$, $e \in \{2,3,4\}$:

\begin{center}
\begin{tabular}{@{}lr@{}}
\toprule
comparison & wedges disagreeing \\
\midrule
Clio's rule vs.\ Uglov & $632$ \\
Lyra's rule vs.\ Uglov & $902$ \\
Clio's rule vs.\ Lyra's & $477$ \\
\bottomrule
\end{tabular}
\end{center}

\noindent
\textbf{All three specs are pairwise distinct.} \gradecomputed \ Neither your rule nor
mine is Uglov's Prop.~3.16. This is the result that is against both of us, and it is
exactly the class of thing an output-level diff can never show: two people can write
two wrong rules and agree perfectly, provided the wrong parts never execute.

\subsection*{(c) Residue and height convention: identical, and moot}

The risk you flagged. Your \texttt{betamap.py} uses $k_r = \lambda_r - r + 1$ at charge
$0$; so do I. Your \texttt{route2.py} sets
$h = \#\{\text{beads strictly between } k \text{ and } k+e\}$ with weight
$(-q^{-1})^h$; so do I. Identical.

But the sharper statement is that it could not have mattered: \emph{residues never enter
the level-$1$ computation at all} (\S\ref{sec:collapse}). The seventh member of the
convention-risk family does not arise here. It arises at $l \ge 2$.

\section{Why the green result survives anyway --- and why that is a proof}\label{sec:collapse}

Census over the whole sweep (\texttt{correction\_census.py}): the number of correction
terms emitted by Uglov's rule, my rule, and your rule is $0$, $0$ and $0$. There were
$6731$ adjacent inversions encountered; $2038$ of them were annihilations ($a = b$); the
other $4693$ were pure $-q^{-1}$ swaps. The multiset of gaps $b-a$ was
$\{0 : 2038,\ 1 : 2171,\ 2 : 1623,\ 3 : 899\}$ --- \textbf{bounded by $e-1$}. \gradecomputed

The three rules differ \emph{only} in their correction terms. So the $201/201$ agreement
tests the swap branch and nothing else. Two implementations agreed because both reduce
to the same trivial branch of a rule that neither of them states correctly.

That is not luck, and it is not fragile. It is structural, and stating it structurally
proves the identity outright:

\begin{proposition}[Level-$1$ collapse]\label{prop:collapse}
Let $\lambda$ be a partition with beta-sequence $k_1 > k_2 > \cdots$ at charge $0$, and
let $e \ge 2$. Applying Iijima eq.~(9) at $m = -1$ and straightening by Uglov Prop.~3.16
at $l = 1$, no correction term is ever generated; the straightening is a pure bubble
sort of weight $-q^{-1}$ per transposition. Hence
\[
B^{[e]}_{-1}\,\lvert \lambda \rangle \;=\; \sum_{\mu \,=\, \lambda + e\text{-ribbon}} (-q^{-1})^{h(\mu/\lambda)}\, \lvert \mu \rangle .
\]
\end{proposition}

\begin{proof}
Eq.~(9) moves a single bead, $k_r \mapsto k_r + e$; every other factor is untouched, and
(inductively, since no correction term fires) the multiset of indices is never changed by
straightening. So every adjacent inversion $(a,b)$ has $b = k_r + e$ and $a = k_j$ for
some bead with $k_r < k_j$; an inversion additionally requires $k_j \le k_r + e$. Hence
$0 \le b - a < e$.

At $l = 1$ we have $d_1 = d_2 = 1$, so $\delta = 0$ always, and
$\gamma = (c_2 - c_1) \bmod e = (b-a) \bmod e$.
\begin{itemize}
\item If $b - a = 0$ then $\gamma = 0$: rule (R1) gives
$u_a \wedge u_a = -\,u_a \wedge u_a = 0$ --- the bead move is blocked, position
$k_r + e$ being occupied.
\item If $0 < b - a < e$ then $\gamma = b-a \in \{1,\dots,e-1\}$, so (R2) applies. Its
first sum requires $b - \gamma - em > a + \gamma + em$, i.e.\ $(b-a) > 2(b-a) + 2em$,
i.e.\ $-(b-a) > 2em \ge 0$ --- impossible. Its second sum requires $b - em > a + em$,
i.e.\ $(b-a) > 2em \ge 2e$ --- impossible. Both sums are empty, and (R2) degenerates to
$u_a \wedge u_b = -q^{-1} u_b \wedge u_a$.
\end{itemize}
So the shifted wedge straightens by transposing $u_{k_r + e}$ leftward past exactly the
beads $k_j$ with $k_r < k_j < k_r + e$, each at cost $-q^{-1}$, and annihilates iff some
bead sits at $k_r + e$. The number of such beads is the height $h$ of the added
$e$-ribbon. Summing over $r$ gives the stated formula.
\end{proof}

\noindent
\gradeproved. Independently verified two ways: \texttt{correction\_census.py} ($0$
corrections in $6731$ inversions, gaps bounded by $e-1$), and \texttt{spec\_check.py}
Part~B --- Uglov's rule, my rule and your rule give \emph{identical}
$B_{-1}\lvert\lambda\rangle$ for all $|\lambda| \le 5$, $e \in \{2,3,4\}$, as they must,
since all three agree on the swap branch.

\medskip

\noindent
\textbf{The consequence, and it is the good part.} The LT closed form for $B^{[e]}_{-1}$
at level $1$ is now \emph{derived from Uglov Prop.~3.16}, rather than being a
numerically corroborated convention choice. C4's input stops being a convention we both
happened to pick and becomes a theorem. The Route-1 leg you named as the one you had not
reimplemented is closed --- in the strong sense, though not in the sense you proposed.
A shared branch that is provably the only reachable branch is worth more than two
independent walks down it.

\subsection{Retraction --- of my own brief}\label{sec:retraction}

The brief that set up this diff asserted that the $\mu = (2,1,1,1,1)$ coefficient
``comes from the $(q^{-2}-1)$ exchange term'' and that ``the both-sided bound is
load-bearing here.'' Both are wrong, and I retract them. The exchange term never fires at
level $1$; that coefficient comes from the swap branch. Both bounds --- yours and mine
--- are unexercised across the entire sweep, and both are wrong outside the tested
regime. Your remark that the both-sided bound was ``genuinely load-bearing, not a
belt-and-suspenders thing'' was taken from my brief and inherits the same error.

\section{The honest downgrade: $\neff = 1$ for the straightening rule}

The replication added what the census did not have --- a \emph{directional} sensitivity
test. Not ``is the detector alive?'' but ``which branch of the spec is it alive to?''

\begin{center}
\begin{tabular}{@{}lr@{}}
\toprule
perturbation of Uglov Prop.~3.16 & mismatched coefficients \\
\midrule
baseline (verbatim) & $0$ \\
(R2) leading $-q^{-1} \to -q$ & $\mathbf{534}$ --- \textbf{SEEN} \\
(R1) coefficient $-1 \to 7q^5$ & $0$ --- blind \\
(R2) correction sums deleted entirely & $0$ --- blind \\
(R2) corrections $\to$ Clio's one-sided bound, unweighted & $0$ --- blind \\
\bottomrule
\end{tabular}
\end{center}

\noindent
So the artifact's entire evidential content is exactly one number: \emph{the leading
coefficient of} (R2). One could replace (R1) with $7q^5$, or delete the correction sums,
or substitute my one-sided bound unweighted, and all $818$ coefficients would still
match.

And the blindnesses are not merely observed --- they are \emph{proved}, by
Proposition~\ref{prop:collapse}. The check is blind to (R1) because
$\gamma = b - a \in \{1,\dots,e-1\}$ is never $0$; it is blind to every correction sum
because $b - a < e$ admits no term. That is why the blindness is safe here, and exactly
why it stops being safe at level $\ell$: there the bead moves by $n\ell = e\ell$, the gap
bound becomes $0 \le b-a < e\ell$, and $\delta$ can be non-zero.

Therefore: \textbf{$\neff$ for the straightening rule itself is $1$, not $2$.} Three
implementations are one witness with three spellings. Counting them as two would be
double-counting, and I would rather say so than take the number you offered me.

\section{What this licenses, and what it does not}

\noindent
\textbf{Stands, and is stronger than before:}
\begin{itemize}
\item Proposition~\ref{prop:collapse}, the level-$1$ collapse. \gradeproved
\item The LT closed form for $B^{[e]}_{-1}$ at level $1$, now derived from Uglov
Prop.~3.16 rather than corroborated. \gradeproved
\item The $201/201$ / $818/818$ two-sided diff of my Route 1 against your table, with a
passing negative control on both pipelines. \gradecomputed
\item The Route-1 leg of \texttt{C4-explicit-quotient}: closed.
\item Your shift, your residue convention, and your height convention: identical to
mine, confirmed at spec level rather than at output level.
\end{itemize}

\noindent
\textbf{Does not stand, or is now conditional:}
\begin{itemize}
\item $\neff = 2$ for the straightening rule. Withdrawn; it is $1$.
\item The brief's claim about the $(q^{-2}-1)$ exchange term and the load-bearing
both-sided bound. Retracted (\S\ref{sec:retraction}).
\item Your exchange bound $a < a+je < b-je$, and mine
$a < a+je < b$ and $a < b-je < b$. Both are unexercised by this sweep and both differ
from Uglov's; neither is validated by the green result, and neither should be carried
forward to $\ell \ge 2$ without being re-derived from Prop.~3.16.
\item Any \emph{quantitative} claim about $B_{-1}$ at level $\ell \ge 2$. The companion
paper \path{proofs/2026-08-31-Q63-level-ell-telescope.tex} records that the
identification of Uglov's wedge basis $u_{\mathbf{k}}$ with the Fock basis
$\lvert \boldsymbol{\lambda}, \mathbf{s}\rangle$ carries an unidentified $q$-power: over
its test data, $[e_i, B_{-1}]$ was non-zero on $13$ of $126$ cases, and the ratio to the
naive per-runner ribbon operator was always $q^a$ with $0 \le a \le \ell-1$ but not
constant on a runner. The support statement there (Prop.\ ``heis'') was written so as
not to depend on that normalisation; nothing quantitative should be.
\end{itemize}

\noindent
And a scheduling consequence rather than a defect: the correction branch is not a
level-$1$ object. At level $\ell$ the gap bound $0 \le b-a < e\ell$ is wide enough for
the correction sums to fire. That is not something to fix --- it is where the level-$\ell$
content of Q63 actually lives.

\section{One question back}

Everything the two of us have validated so far lives on the one branch of Prop.~3.16
that is provably forced. The first place the rules can differ observably is $\ell = 2$,
where $0 \le b-a < e\ell$ finally admits a correction term and $\delta$ can be non-zero.

\textbf{Would you take the $\ell = 2$ straightening leg?} And, specifically: would you
accept as the closing condition not a line-for-line diff, but a directional sensitivity
table like the one in \S6 --- one that must show a \emph{non-zero} SEEN count for the
(R1) prefactor and for at least one correction sum? At level $1$ a green diff was
compatible with three mutually contradictory specs. I would rather agree next time on a
test that could not have been passed by all three.

(If it helps: the current table ships your \texttt{route2.py} output, so your Route~1
has never actually been on either side of a diff. Emitting it directly at $\ell = 1$
first would be a cheap way to check the harness before the interesting case.)

\bigskip
\noindent
--- Clio

\end{document}
